% !TEX program = lualatex
% ---------------------------------------------------------------------------
%  日本語論文（二段組）
%  LuaLaTeX + jlreq（日本語組版処理の要件に準拠）。和文は原ノ味フォント、
%  欧文は Latin Modern。参考文献は BibTeX（refs.bib）。
%  LuaLaTeX と原ノ味フォントは初回コンパイル時にダウンロードして追加する（約 32 MB）。
%  研究内容・数値・著者はすべて例なので、自分の内容に置き換えて使う。
% ---------------------------------------------------------------------------
\documentclass[paper=a4paper,fontsize=10pt,twocolumn]{jlreq}

\usepackage[haranoaji]{luatexja-preset} % 原ノ味フォント（TeX Live 同梱）
\usepackage{amsmath,amssymb}
\usepackage{graphicx}
\usepackage{booktabs}
\usepackage[hidelinks]{hyperref}

\title{深層学習を用いた手書き数式認識の高精度化に関する研究}
\author{山田 太郎\thanks{〇〇大学大学院 情報理工学研究科} \and 佐藤 花子\thanks{同上}}
\date{2026年10月27日}

\begin{document}

\twocolumn[
  \maketitle
  \begin{abstract}
    本研究では，手書き数式画像から \LaTeX{} 記法を生成する手法を提案する．
    従来手法では長い数式における構造の誤認識が課題であった．
    そこで，木構造を考慮したデコーダを導入し，認識精度を改善した．
    公開データセットを用いた評価実験の結果，提案手法は既存手法と比較して
    完全一致率を 4.2 ポイント向上させた．
  \end{abstract}
  \noindent\textbf{キーワード：}手書き数式認識，深層学習，系列変換
  \vspace{1em}
]

\section{はじめに}
数式は科学技術文書において不可欠な要素である．
しかし，手書き数式をデジタル化する作業は依然として負担が大きい\cite{zhang2017}．
本稿の貢献は以下の通りである．
\begin{itemize}
  \item 木構造を考慮したデコーダの提案
  \item 大規模データセットにおける有効性の検証
\end{itemize}
なお，本稿の組版には LuaLaTeX と jlreq を用いた\cite{okumura2023}．

\section{関連研究}
系列変換モデルに基づく手法として，注意機構を用いたモデルが広く用いられている\cite{vaswani2017}．

\section{提案手法}
入力画像 $\mathbf{X}$ に対し，出力系列 $\mathbf{y} = (y_1, \dots, y_T)$ の条件付き確率を
式~\eqref{eq:prob} のようにモデル化する．
\begin{equation}
  p(\mathbf{y} \mid \mathbf{X}) = \prod_{t=1}^{T} p(y_t \mid y_{<t}, \mathbf{X})
  \label{eq:prob}
\end{equation}

\begin{figure}[tb]
  \centering
  \fbox{\parbox[c][3cm][c]{0.8\linewidth}{\centering 図をここに挿入}}
  % \includegraphics[width=0.9\linewidth]{figure.pdf}
  \caption{提案手法の概要}
  \label{fig:overview}
\end{figure}

全体の構成を図~\ref{fig:overview} に示す．

\section{評価実験}
実験結果を表~\ref{tab:result} に示す．

\begin{table}[tb]
  \centering
  \caption{完全一致率の比較}
  \label{tab:result}
  \begin{tabular}{lc}
    \toprule
    手法 & 完全一致率 (\%) \\
    \midrule
    従来手法 A & 52.1 \\
    従来手法 B & 54.8 \\
    提案手法   & \textbf{59.0} \\
    \bottomrule
  \end{tabular}
\end{table}

\section{おわりに}
本稿では，手書き数式認識のための新たな手法を提案した．
今後の課題として，複数行にわたる数式への拡張が挙げられる．

\section*{謝辞}
本研究は JSPS 科研費 JP00000000 の助成を受けたものである．

\bibliographystyle{unsrt}
\bibliography{refs}

\end{document}
